\documentclass{article}
\usepackage{pxLST}
\begin{document}
% Terminal content is verbatim; dollar and quotes remain literal.
\begin{pxTerminal}{Shell session}
$ printf '%s\n' "$HOME" # literal shell prompt
/home/student
\end{pxTerminal}

% The same file renderer handles a full file and a selected range.
\pxCodeFile[pxJava]{sample.java}[Java source]
\pxCodeChunk[pxJava]{sample.java}{2}{5}{Java excerpt}

% Retain the original compact/annotated environments and their behavior.
\begin{miniCode}
int a = 2; // literal # % and $
\end{miniCode}
\begin{CodeBox}{new-label}{Code with note}{note}{Annotation}
int b = 3;
\end{CodeBox}
\begin{CodeBox*}{new-break}{Breakable code}{note}{Annotation}
int c = 4;
\end{CodeBox*}
% Original public contracts must remain operational.
\begin{miniCode}
int old = 1;
\end{miniCode}
\begin{CodeBox}{old-label}{Legacy box}{note}{Note}
int oldBox = 2;
\end{CodeBox}
\begin{CodeBox*}{old-break}{Legacy break}{note}{Note}
int oldBreak = 3;
\end{CodeBox*}
\pxCodeIn[pxJava]{int oldCall = 4;}
New inline: \pxCodeInline[pxJava]|int newCall = 5;|.
\typeout{PXLST-RECOVERY-PASS}
\end{document}
